% ============================================================
% +10分·阅读与写作（杂志页码 46–53）
% 思维表达 · 文艺鉴赏 · 经典洞见
% ============================================================
\colchapter{+10分·阅读与写作}

\section*{思维表达：具备“松弛感”的人，有选择“不松弛”的自由}
\auth{编写者：话秋}

2022年，互联网上兴起了一个词——“松弛感”。它源于一次乘机经历：一位网友发帖称，自己一家人出门旅游，因孩子证件过期，妈妈只能留下来陪孩子，所有人的行李都被退回了。这本是一件让人焦头烂额的事，没想到这家人完全没有吵闹、生气，而是平静地重新安排行程，在电话那头语气轻松地处理后续。于是，帖主感叹道：“原来这个世界上真的有‘松弛感’这种生活状态。”

“松弛感”就此走红，很快成为人们对理想生活状态的想象：不纠结、不拧巴，遇事不慌、从容不迫。有人羡慕，有人努力模仿，也有人反问道：难道“松弛”真的可以学来吗？

其实，“松弛感”并不是一种可以刻意表演的姿态，而是一种内在的从容。真正具备“松弛感”的人，往往并不是没有压力、没有目标，而是他们清楚地知道自己要什么、能接受什么，因而不会在无谓的事情上耗费心力。换言之，松弛的背后，是“想得明白”的底气。

更进一步说，具备“松弛感”的人，恰恰拥有选择“不松弛”的自由。真正松弛的人，可以为热爱的事业全力以赴、为重要的事情紧张筹备，因为这份“紧绷”是他们主动选择的；而假装松弛的人，却可能在逃避责任与拖延中失去了自己的节奏。松弛与否，关键在于是否掌握着人生的主动权。

\begin{quoteBox}{哲思笔记}
“松弛感”不等于“躺平”，更不是对生活的敷衍。它是一种经历过磨砺之后形成的通透：懂得抓大放小，懂得接纳不完美，也懂得在全力以赴与顺其自然之间找到平衡。松弛的底色，是看清生活真相之后依然热爱生活。
\end{quoteBox}

\clearpage
\section*{没有“松弛感”也没关系}
\auth{文/赵小小月}

“松弛感”这个词，最近在网络上很火。它描述的是一种从容不迫、不慌不忙的生活状态：遇到突发状况不急不躁，面对未知未来不忧不惧，仿佛一切尽在掌握之中。

可是，当“松弛感”被当作一种标准去衡量自己时，很多人反而更焦虑了——为什么别人能那么松弛，而我总是紧绷着？为什么我学不会松弛？这样的追问，让“松弛感”从一个美好的向往，变成了一种新的心理负担。

其实，没有“松弛感”也没关系。每个人的性格、经历、处境各不相同。有人天生随和，有人习惯思虑周全；有人能在短时间内调整好心态，有人则需要更长的时间消化情绪。紧绷，未必是弱点，它可能只是你对生活认真的一种方式。过分要求自己“松弛”，恰恰是一种新的“紧绷”。

更重要的是，真正的“松弛感”不是表演给别人看的，而是源于内心的安全感和掌控感。它建立在对自我的接纳之上：接纳自己的情绪，接纳自己的局限，也接纳生活的不完美。当我们不再用“松弛感”来苛责自己，反而更容易在日复一日的努力中，找到属于自己的节奏。

所以，如果你暂时没有“松弛感”，不必自责。允许自己紧张，允许自己焦虑，允许自己慢慢来。生活的答案，从来不只有一个标准。你不需要活成别人眼中的“松弛”，你只需要活成自己舒服的样子。

\begin{quoteBox}{写作分析}
本文从“松弛感”成为网络热词谈起，先呈现现象，再指出“以松弛为标尺反致焦虑”的矛盾，接着给出观点——没有松弛感也没关系，最后落点到“接纳自我、找到自己的节奏”。层层递进，逻辑清晰，为热点类议论文提供了范例。
\end{quoteBox}

\clearpage
\section*{从“躺平”到“松弛感”，拒绝“内卷”的年轻人到底想要什么？}
\auth{文/之初}

从“躺平”到“内卷”，再到“松弛感”，近几年的互联网热词，记录着年轻人的心态变化。

“内卷”是同一赛道上的过度竞争：别人多学一小时，我也要多学一小时；别人加班到九点，我就加班到十点。在这样的竞争里，投入越来越多，收获却未必同步增长，人也被裹挟着不断加速。“躺平”则是对“内卷”的极端反弹：既然怎么努力都追不上，那就干脆不努力了。然而，真正“躺平”的人并不多，大多数人只是嘴上说说，身体依然诚实地卷着。

“松弛感”的流行，则像是年轻人在“内卷”与“躺平”之间找到的第三条路：不盲目追逐，也不彻底放弃；有自己的节奏，也接纳生活的意外。比起“赢了所有人”，“松弛感”更看重“与自己和解”。

那么，拒绝“内卷”的年轻人到底想要什么？他们想要的，从来不是不劳而获，而是付出能被看见、努力能有意义；他们想要的，不是彻底的躺平，而是生活不只工作、人生不只KPI；他们想要的，是属于自己的时间，是选择的权利，是“即使慢一点也没关系”的从容。

可以说，“松弛感”的背后，是一种更加成熟的价值观：人生不是一场非赢即输的比赛，而是一段可以按自己节奏行走的旅程。与其在别人的赛道里精疲力竭，不如在自己的旷野上从容奔跑。

\begin{quoteBox}{考场范本旁栏}
写作时可将“内卷—躺平—松弛感”作为一组动态概念进行梳理，用时间线或对比结构展现青年心态的变迁。注意避免简单褒贬，要看到三者之间的内在关联，才能在立意上胜人一筹。
\end{quoteBox}

\clearpage
\section*{名师小书卷：能够飞翔的人生}
\auth{考场范本 · 上海市七宝中学 俞佳（指导老师：张君平）}

人生而渴望飞翔。然而，能够飞翔的人生，从来不是与生俱来的天赋，而是一步步积累、一次次突破的馈赠。

能够飞翔的人生，始于热爱的托举。鸟的翅膀，生来就为飞翔；人的翅膀，却需要在热爱中找到。热爱，让我们愿意为之付出时间与心血，让我们在枯燥的训练中依然甘之如饴。就像那些在各自领域翱翔的“飞人”们，无论是赛场上的运动员，还是舞台上的艺术家，他们之所以能飞得高、飞得远，无不是因为心中有一团不灭的热爱之火。

能够飞翔的人生，成于坚持的磨砺。没有哪一种飞翔是一帆风顺的。雏鹰学飞，要经历无数次的跌落；人类圆梦，要跨越数不清的沟坎。坚持，让一次次看似微小的努力累积成质变，让一次次失败成为向上的阶梯。正如那句名言所说：“每一个优秀的人，都有一段沉默的时光，那段时光，是付出了很多努力，却得不到结果的日子，我们把它叫作扎根。”

能够飞翔的人生，贵在方向的笃定。飞得再高，若方向错了，也只会离目标越来越远。真正的飞翔者，懂得在纷繁的诱惑中守住初心，在嘈杂的声音里听见自己。他们知道，飞翔的意义不在于飞得多高，而在于是否飞向了心中的远方。

能够飞翔的人生，不是完美的人生，而是完整的人生。它接纳笨拙的起步，包容踉跄的跌落，也珍藏每一次逆风向上的瞬间。愿我们都能在自己的天空里，学会飞翔，且行且歌。

\begin{quoteBox}{技法点击}
本文采用“总—分—总”结构，以“热爱—坚持—方向”三个分论点展开，结尾回扣主题并升华。语言典雅而有节奏，善用比喻与排比，是典型的考场佳作写法。
\end{quoteBox}

\clearpage
\section*{文艺鉴赏：《珠峰队长》：逆流而上，攀登永不止步}
\auth{编写者：庆嫣}

2022年，登山纪录片《珠峰队长》上映。影片跟随一支由普通人组成的登山队，历时四十余天，记录了他们攀登珠穆朗玛峰的全过程。没有剧本，没有演员，只有真实的风雪、真实的疲惫和真实的抉择——而这，恰恰是它最动人的地方。

片中的“珠峰队长”苏拉王平，出生在四川阿坝州的一个小山村。1998年，一次偶然的机会，他成为高山向导。此后二十多年，他无数次往返于雪山之巅，把一批又一批登山者安全地带上去、带回来。为了拍一部属于中国登山者的纪录片，他带着队伍和摄影团队，在海拔八千米以上的“死亡地带”，完成了人类历史上首次无人机全程航拍珠峰的壮举。

攀登的意义是什么？影片给出了朴素而动人的回答：对苏拉王平来说，登山是谋生的职业，更是热爱的事业；对队伍里的普通登山者来说，登上峰顶，是给自己的一个交代，是“我也可以”的证明。攀登的过程中，有高原反应的折磨，有恶劣天气的考验，有放弃的念头一闪而过，但最终，他们都选择了继续向上。

《珠峰队长》最动人的，不是登顶那一刻的欢呼，而是那些“逆流而上”的瞬间：是在疲惫到极点时依然迈出的下一步，是队友之间相互扶持的一句“加油”，是面对危险时依然选择信任的勇气。人生亦如登山，每个人都在攀登属于自己的那座山。重要的不是登顶的高度，而是向上攀登的每一步。

\begin{quoteBox}{人物档案}
苏拉王平，四川阿坝人，登山向导、纪录片导演。自1998年起从事高山向导工作，累计登顶珠峰多次，2019年率队完成珠峰全程航拍纪录片的拍摄。他说：“山不会向你走来，只有你向山走去。”
\end{quoteBox}

\clearpage
\section*{观影评说：始于梦想，忠于内心}
\auth{文/九九}

看完纪录片《珠峰队长》，我记住的不是峰顶的风景，而是出发时那一张张平凡的脸。

影片中的登山队员，有开火锅店的老板，有做建材生意的商人，有朝九晚五的上班族……他们来自各行各业，年龄各异，却因为同一个梦想聚在一起。为了这次攀登，他们进行了长期的高原适应训练，付出了时间、金钱和汗水。有人问：花这么多钱、冒这么大风险，值得吗？影片里没有直接回答，但每一个镜头都在回应：因为热爱，所以值得。

“始于梦想”，是他们出发的理由；“忠于内心”，是他们坚持的动力。在海拔八千米的雪线上，支撑他们走下去的，不是外在的掌声，而是内心深处那个“我想试一试”的声音。登顶之后，有人激动落泪，有人平静微笑，但所有人都明白：这一路的经历，远比山顶的风景更加珍贵。

这让我想到，人生中的许多事，其实都像一次攀登。我们未必都能登顶，但那些为目标全力以赴的日子，那些与同伴并肩前行的时刻，都会成为生命里最闪亮的记忆。梦想或许遥远，但只要开始行动，就已经站在了通往它的路上。

\begin{quoteBox}{网友评论}
@高山流水：这才是真正的“高燃”！不是特效，是真实的风雪和真实的勇气。
@云端漫步：看哭了。平凡人的梦想，最是动人。
@追光者：人生就是一场攀登，重要的不是登顶，而是向上。
\end{quoteBox}

\clearpage
\section*{名师小书卷：攀登，前行者生命的恒久状态}
\auth{考场范本 · 河北省秦皇岛市北戴河中学 靳国申}

“攀登”一词，不仅指向高山，更指向一种生命的姿态。对前行者而言，攀登是恒久的状态——无论身处何种境地，都保持向上的姿态，都怀有前行的渴望。

攀登，是困境中的突围。被誉为“中国治沙第一人”的董鸿儒，几十年如一日扎根荒山，用一株株树苗向风沙宣战，把不毛之地变成绿色屏障。他攀登的，是一座没有终点的“山”。15岁的也门少女亨德·扎扎，在战火纷飞中坚持乒乓球训练，哪怕没有正规场地，也要在瓦砾间挥拍。她攀登的，是一条充满荆棘的路。这些前行者告诉我们：困境不是止步的理由，而是攀登的起点。

攀登，是平凡中的超越。运动员凯雷日复一日地在跑道上奔跑，用汗水浇灌出属于普通人的冠军梦；更多不知名的劳动者，在各自的岗位上日复一日地精进，把平凡的日子过成闪光的岁月。攀登，未必是惊天动地的壮举，也可以是每一次“再努力一点”的坚持。

攀登，更是信念的坚守。没有信念的攀登，只是盲目的攀爬；有了信念的支撑，攀登便有了方向与力量。那些在历史长河中留下印记的前行者，无一不是以信念为杖，以勇气为履，在属于自己的山路上，一步一个脚印地向上。

人生的山，没有顶峰。但正因如此，攀登才成为前行者生命恒久的状态——永远向上，永远向前，永远在路上。

\clearpage
\section*{课堂演练 + 考场范本：执着攀登，勇毅前行}
\auth{考场范本 · 李冉}

【课堂演练材料】“没有比人更高的山，没有比脚更长的路。”攀登，是人类挑战自我、超越极限的永恒主题。从珠峰之巅的国旗，到航天员漫步太空；从万佐成夫妇十八年如一日的“抗癌厨房”，到袁隆平一生守望稻田的“禾下乘凉梦”，无数中华儿女以攀登者的姿态，书写着奋斗者的答卷。

【考场范本】

人生如登山，山高路远，唯有执着攀登、勇毅前行，方能领略顶峰的风景。

执着攀登，需要“不破楼兰终不还”的决心。目标一旦确立，便义无反顾。万佐成夫妇在肿瘤医院旁支起“抗癌厨房”，一干就是十八年。灶火不熄，是因为他们认定了“要让病人吃上一口家里的饭”；初心不改，是因为他们知道，这一方小小的厨房，承载着多少家庭的牵挂。这平凡的坚守，何尝不是一种攀登？

执着攀登，需要“越是艰险越向前”的勇气。攀登路上，风雪常在。袁隆平一生与水稻为伴，从“三系法”到“两系法”，从亩产700公斤到亩产1500公斤，每一次突破，都是在质疑与困难中闯出来的。航天人更是如此，从东方红一号到载人航天，一代代航天人用勇气和智慧，把“不可能”变成“一定能”。

执着攀登，需要“功成不必在我”的胸襟。很多攀登者，终其一生也未必能登上峰顶，但他们的每一步，都为后来者铺就了道路。正是无数无名者的默默攀登，才汇成了民族复兴的壮阔征程。

山高人为峰。让我们以攀登者的姿态，在人生的征途上执着前行，勇毅向前，用奋斗书写属于这个时代的答卷。

\begin{quoteBox}{技法点击}
考场范本紧扣“执着攀登”主题，以“决心—勇气—胸襟”三个维度展开论证，层次分明；每个分论点均结合时代人物事例，内容充实；结尾点题升华，气势昂扬。
\end{quoteBox}
